\documentclass[12pt,a4paper]{article}
\usepackage[utf8]{inputenc}
\usepackage{amsmath,amssymb,amsthm}
\usepackage{algorithm}
\usepackage{algorithmic}
\usepackage{geometry}
\geometry{margin=1in}

\newtheorem{theorem}{Theorem}
\newtheorem{lemma}[theorem]{Lemma}
\newtheorem{corollary}[theorem]{Corollary}
\newtheorem{proposition}[theorem]{Proposition}
\theoremstyle{definition}
\newtheorem{definition}{Definition}
\theoremstyle{remark}
\newtheorem*{remark}{Remark}

\title{Project Euler Problem 23: Non-Abundant Sums}
\author{}
\date{}

\begin{document}
\maketitle

\section{Problem Statement}

A positive integer $n$ is \emph{abundant} if the sum of its proper divisors exceeds~$n$. Find the sum of all positive integers that cannot be written as the sum of two abundant numbers, given that every integer greater than $28123$ admits such a representation.

\section{Definitions}

\begin{definition}[Abundant Number]
A positive integer $n$ is \emph{abundant} if $s(n) > n$, where $s(n) = \sigma_1(n) - n$.
\end{definition}

\begin{definition}[Abundant-Sum Representability]
A positive integer $m$ is \emph{abundant-sum representable} if there exist abundant numbers $a, b$ with $m = a + b$. Denote the set of such integers by $\mathcal{A}_2$.
\end{definition}

\section{Theorems and Proofs}

\begin{theorem}[Smallest Abundant Number]
The smallest abundant number is $12$.
\end{theorem}

\begin{proof}
Exhaustive verification shows $s(n) \leq n$ for all $n \leq 11$ (with $s(6) = 6$ being perfect). For $n = 12$: $s(12) = 1 + 2 + 3 + 4 + 6 = 16 > 12$.
\end{proof}

\begin{theorem}[Closure Under Scaling]
If $n$ is abundant and $k \geq 2$, then $kn$ is abundant.
\end{theorem}

\begin{proof}
For each $d \mid n$, the integer $kd$ divides $kn$, so $\{kd : d \mid n\} \subseteq \{d : d \mid kn\}$. Hence
\[
\sigma_1(kn) \geq \sum_{d \mid n} kd = k\,\sigma_1(n) > 2kn,
\]
where the last inequality uses $\sigma_1(n) > 2n$. Thus $s(kn) > kn$.
\end{proof}

\begin{corollary}
Every positive multiple of $12$ that is at least $24$ is abundant.
\end{corollary}

\begin{theorem}[Finite Complement of $\mathcal{A}_2$]
Every integer greater than $28123$ belongs to $\mathcal{A}_2$.
\end{theorem}

\begin{proof}
By the preceding corollary, all multiples of $12$ from $24$ onward are abundant. Since abundant numbers have positive asymptotic density (Davenport, 1933: density $\approx 0.2477$), the sumset $\mathcal{A}_2$ contains all sufficiently large integers. The bound $28123$ has been verified computationally.
\end{proof}

\begin{remark}
The largest integer not in $\mathcal{A}_2$ is $20161$, so $28123$ is a conservative upper bound.
\end{remark}

\section{Algorithm}

\begin{algorithm}[H]
\caption{Sum of Non-Abundant Sums}
\begin{algorithmic}[1]
\REQUIRE Upper bound $N = 28123$
\ENSURE $\displaystyle\sum_{\substack{n=1 \\ n \notin \mathcal{A}_2}}^{N} n$
\STATE Sieve: for $i = 1$ to $N$, add $i$ to $\mathrm{spd}[j]$ for each multiple $j = 2i, 3i, \ldots \leq N$
\STATE $A \gets \{i \in [1, N] : \mathrm{spd}[i] > i\}$
\STATE Initialize $\mathrm{repr}[1..N] \gets \mathrm{false}$
\FOR{$i = 1$ \TO $|A|$}
    \FOR{$j = i$ \TO $|A|$}
        \IF{$A[i] + A[j] > N$}
            \STATE \textbf{break}
        \ENDIF
        \STATE $\mathrm{repr}[A[i] + A[j]] \gets \mathrm{true}$
    \ENDFOR
\ENDFOR
\RETURN $\sum_{n=1}^{N} n \cdot [\neg\,\mathrm{repr}[n]]$
\end{algorithmic}
\end{algorithm}

\section{Complexity Analysis}

\begin{proposition}
The algorithm runs in $O(N \log N + A^2)$ time and $O(N)$ space, where $A \approx 0.2477\,N$.
\end{proposition}

\begin{proof}
The sieve in Step~1 performs $\sum_{i=1}^{N}(\lfloor N/i\rfloor - 1) = N H_N - N = O(N \log N)$ operations. The double loop over abundant numbers executes at most $\binom{A}{2} + A = O(A^2)$ iterations. The boolean and divisor-sum arrays each require $O(N)$ space.
\end{proof}

\section{Answer}

\[
\boxed{4179871}
\]

\end{document}
